\section{Unified core proof chain}

This chapter expands the unified core proof of documents 36.30.22 and
36.30.23.

\subsection{Candidate parameterisation}

The block-head state is
\begin{equation*}
  r=\frac{A_j+5^jq}{2^{\wordWeight_j}}.
\end{equation*}
The predecessor candidate is
\begin{equation*}
  x=2^{e-1}g+\delta5^{j-1},
\qquad g=\Orbit_7(j-1).
\end{equation*}

Status:
\begin{itemize}[leftmargin=*]
\item Math: proven.
\item Lean: \texttt{reset\_head\_predecessor},
  \texttt{candidateX\_of\_reset\_and\_terminal},
  \texttt{first\_block\_terminal\_eq} compiled.
\end{itemize}

\subsection{Reset-head predecessor}

\begin{lemma}
If $1\le j$, the reset equation
$\mathrm{ResetHeadEq}(s,j,k,t,\delta,r_j)$ holds, and
\begin{equation*}
  r_j=\frac{5x+1}{2^t},
\end{equation*}
then
\begin{equation*}
  x=5^k\,s+\delta\,5^{j-1}-1.
\end{equation*}
This is \texttt{reset\_head\_predecessor}.
\end{lemma}

\subsection{First-block terminal equation}

\begin{lemma}
If $5g_{\mathrm{prev}}+1=2^eg$ and
\begin{equation*}
  r=\frac{5g_{\mathrm{prev}}+1}{2},
\end{equation*}
then
\begin{equation*}
  r=2^{e-1}g.
\end{equation*}
With $k=0$ and $s=2^{e-1}g+1$, the reset predecessor becomes
\begin{equation*}
  x=\texttt{candidateX}\ j\ e\ g\ \delta.
\end{equation*}
The Lean names are \texttt{first\_block\_terminal\_eq} and
\texttt{candidateX\_of\_reset\_and\_terminal}.
\end{lemma}

\subsection{Reset q0 form}

\begin{lemma}
From the exact identity
\begin{equation*}
  A_j+5^jq=2^L(B+\delta5^j),
\end{equation*}
with $0<B<5^j$ and $A_j<5^j$, there exists $m<2^L$ with
\begin{equation*}
  q=m+\delta\,2^L.
\end{equation*}
This is \texttt{reset\_q0\_form}.
\end{lemma}

\subsection{Block-head identity}

\begin{lemma}
Under the block-head representation and the reset equation,
\begin{equation*}
  A_j+5^jq
    =2^{W_p}\left(5^{k+1}s-4+\delta5^j\right).
\end{equation*}
This is \texttt{block\_head\_identity\_of\_reset}.
\end{lemma}

\subsection{Orbit segment word}

For $d$ steps starting at depth $j-1$, the segment word satisfies
\begin{equation*}
  2^{W}x=5^d\,g+\wordA(w),
\end{equation*}
where $W$ is the total segment weight.  This is
\texttt{orbitSegmentWord\_equation}.  The orbit-level exclusions
\texttt{d1\_exclusion\_of\_orbit} and
\texttt{d2\_exclusion\_of\_orbit} are compiled.

\subsection{Derivation of the segment equation}

Let the segment word $w$ have length $d$, total weight $W=W_d$, and
prefix weights $W_0,\dots,W_d$.  The exact word-orbit identity applied
to the segment from $g$ to $x$ gives
\begin{equation*}
  2^W x=5^d\,g+\wordA(w),
\end{equation*}
with
\begin{equation*}
  \wordA(w)=\sum_{i=0}^{d-1}2^{W_i}\,5^{d-1-i}.
\end{equation*}
Substituting the candidate
\begin{equation*}
  x=2^{e-1}g+\delta5^{j-1}
\end{equation*}
and expanding gives the exact segment equation
\begin{equation*}
  2^{W+e-1}g+2^W\delta5^{j-1}
    =5^d\,g+\sum_{i=0}^{d-1}2^{W_i}5^{d-1-i}.
\end{equation*}
This is the candidate form
\texttt{orbitSegmentWord\_\allowbreak candidate\_\allowbreak equation}.
No approximation enters: the equation is the word-orbit identity with
the parameterisation substituted.

\subsection{Size bounds used in the exclusions}

The exclusion arguments use two size bounds on $g=\Orbit_7(j-1)$:
\begin{equation*}
  g<\frac{5^{j-1}}{4},
  \qquad
  g<\frac{5^{j-1}}{2^{e-1}}.
\end{equation*}
The first is used by the $d=1$ exclusion and by the $e\ge3$ cases;
the second, with $e\ge2$, is used by the $d=2$ size analysis.  Both
are consequences of the no-$H_{\ge}$ premises together with the exact
orbit bound; in the Lean statements they are carried as hypotheses.
The paper records them as premises rather than re-deriving them at
each exclusion.

\subsection{Segment exclusions}

The exact segment equation is
\begin{equation*}
  2^{W+e-1}g+2^W\delta5^{j-1}
    =5^d\,g+\sum_{i=0}^{d-1}2^{W_i}5^{d-1-i}.
\end{equation*}

For small $d$, the segment word molecule and total weight have the
explicit forms:
\begin{center}
\small
\begin{tabular}{@{}lll@{}}
\toprule
$d$ & $\wordA(w)$ & $W$ \\
\midrule
1 & $1$ & $u_1$ (in the $d=1$ branch, $u_1=1+4a$) \\
2 & $5+2^{u_1}$ & $u_1+u_2$ \\
3 & $25+5\cdot2^{u_1}+2^{u_1+u_2}$ & $u_1+u_2+u_3$ \\
\bottomrule
\end{tabular}
\end{center}

For the surviving branches the equations reduce as follows.  With
$\delta=1$, $e=2$, $u_1=1$ the $d=2$ equation becomes
\begin{equation*}
  17g=4\cdot5^{j-1}-7.
\end{equation*}
With $u_1=u_2=1$, $a=0$ (so $u_3=1$) the $d=3$ equation becomes
\begin{equation*}
  16g+8\cdot5^{j-1}=125g+39,
\end{equation*}
hence
\begin{equation*}
  g=\frac{8\cdot5^{j-1}-39}{109}.
\end{equation*}

The exclusions are:
\begin{enumerate}[leftmargin=*]
\item $d=1$: impossible by the size bound $g<5^{j-1}/4$.
\item $d=2$: modular contradiction modulo $16$.
\item $d=3$: first big step at depth $j+4$ with weight $6$.
\item $d\ge4$: first big step weight at least $5$, or the
  $e=3,j=17$ modular exclusion.
\end{enumerate}

Status: all four exclusions are math-proven and Lean-compiled.

\subsection{Worked $d=1$ exclusion}

The $d=1$ segment equation is
\begin{equation*}
  5g+1=2^{1+4a}\left(2^{e-1}g+\delta5^{j-1}\right),
\end{equation*}
i.e.
\begin{equation*}
  5g+1=2^{e+4a}g+\delta2^{1+4a}5^{j-1}.
\end{equation*}
If $e\ge3$ or $a\ge1$, then $2^{e+4a}\ge8$, and the second term is
positive, so the right-hand side exceeds $5g+1$, contradicting the
equality.

If $e=2$ and $a=0$, the equation reduces to
\begin{equation*}
  5g+1=2\left(2g+\delta5^{j-1}\right),
\end{equation*}
so
\begin{equation*}
  g+1=2\delta5^{j-1}.
\end{equation*}
Since $\delta\ge1$, this gives $g\ge2\cdot5^{j-1}-1$, contradicting
the size bound $g<5^{j-1}/4$.

\subsection{Full $d=2$ size analysis}

The $d=2$ segment equation is
\begin{equation*}
  \left(2^{u_1+e}-25\right)g
    +2^{u_1+1}\delta5^{j-1}=5+2^{u_1}.
\end{equation*}
The right-hand side is at most $9$ for $u_1\in\{1,2\}$, so the left-hand
side must be small.  In particular
\begin{equation*}
  2^{u_1+e}<25,
\end{equation*}
which leaves only
\begin{equation*}
  (u_1,e)\in\{(1,2),(1,3),(2,2)\}.
\end{equation*}

For $u_1=1,e=2$, the equation becomes
\begin{equation*}
  17g=4\delta5^{j-1}-7.
\end{equation*}
The size bound is $g<5^{j-1}/2$.  For $\delta=3$ the right-hand side
is about $12\cdot5^{j-1}/17$, which exceeds the bound; only
$\delta=1$ survives.

For $u_1=1,e=3$, the equation becomes
\begin{equation*}
  9g=8\delta5^{j-1}-7,
\end{equation*}
with the stronger size bound $g<5^{j-1}/4$.  Even $\delta=1$ gives
$g\approx8\cdot5^{j-1}/9$, contradicting the bound.

For $u_1=2,e=2$, the equation becomes
\begin{equation*}
  9g=8\delta5^{j-1}-9,
\end{equation*}
with the size bound $g<5^{j-1}/2$; again $\delta=1$ already exceeds it.

Hence the only surviving branch is $\delta=1$, $e=2$, $u_1=1$.  The
subsequent modular contradiction is given in the next subsection.

\subsection{Worked modular computation for $d=2$}

For $d=2$ the surviving branch is $\delta=1$, $e=2$, $u_1=1$, which
reduces the segment equation to
\begin{equation*}
  17g=4\cdot5^{j-1}-7.
\end{equation*}
Modulo $17$ this gives
\begin{equation*}
  4\cdot5^{j-1}\equiv7,
\end{equation*}
and since $4^{-1}\equiv13\pmod{17}$,
\begin{equation*}
  5^{j-1}\equiv7\cdot13\equiv6\pmod{17}.
\end{equation*}
The powers of $5$ modulo $17$ have order $16$, and
\begin{equation*}
  5^3=125\equiv6\pmod{17},
\end{equation*}
so $j-1\equiv3\pmod{16}$.

The candidate congruence $x\equiv743\pmod{1280}$ with
$x=2g+5^{j-1}$ then forces
\begin{equation*}
  5^{j-1}\equiv45\pmod{256}.
\end{equation*}
Since $5^7=78125\equiv45\pmod{256}$ and the order of $5$ modulo $256$
is $64$, this gives $j-1\equiv7\pmod{64}$.  The two congruences
contradict each other modulo $16$.

\subsection{Residue ledger}

\begin{center}
\small
\begin{tabular}{@{}p{0.10\textwidth}p{0.56\textwidth}p{0.28\textwidth}@{}}
\toprule
Case & Congruence and role & Lean \\
\midrule
$d=1$ & no congruence; the size bound $g<5^{j-1}/4$ suffices &
  \texttt{d1\_\allowbreak exclusion} \\
$d=2$ & $5^{j-1}\equiv6\pmod{17}$ and $x\equiv743\pmod{1280}$ force
  $j-1\equiv3\pmod{16}$ and $5^{j-1}\equiv45\pmod{256}$ &
  \texttt{d2\_\allowbreak exclusion} \\
$d=3$ & $j-1\equiv923\pmod{1728}$ fixes the unique surviving family &
  \texttt{d3\_\allowbreak family\_\allowbreak big\_\allowbreak
  weight\_\allowbreak excluded} \\
$d\ge4$ & the $e=3$, $j=17$ branch is excluded by residues modulo
  $640$ or $1280$ &
  \texttt{dge4\_\allowbreak e3\_\allowbreak j17\_\allowbreak t1\_\allowbreak
  excluded}, \texttt{dge4\_\allowbreak e3\_\allowbreak j17\_\allowbreak
  t2\_\allowbreak delta3\_\allowbreak excluded} \\
\bottomrule
\end{tabular}
\end{center}

\subsection{Full $d=3$ survivor derivation}

The surviving $d=3$ parameters are
\begin{equation*}
  t_j=2,\qquad \delta=1,\qquad e=2,\qquad u_1=u_2=1,\qquad a=0,
\end{equation*}
which forces $u_3=1$.  The segment word molecule is
\begin{equation*}
  \wordA(w)=25+5\cdot2^{u_1}+2^{u_1+u_2}
    =25+10+4=39,
\end{equation*}
and the total segment weight is
\begin{equation*}
  W=u_1+u_2+u_3=3.
\end{equation*}
The segment equation
\begin{equation*}
  2^{W+e-1}g+2^W\delta5^{j-1}
    =5^d\,g+\wordA(w)
\end{equation*}
therefore reads
\begin{equation*}
  2^{3+1}g+2^3\cdot5^{j-1}
    =5^3g+39,
\end{equation*}
i.e.
\begin{equation*}
  16g+8\cdot5^{j-1}=125g+39,
\end{equation*}
so
\begin{equation*}
  109g=8\cdot5^{j-1}-39.
\end{equation*}

The integrality of $g$ requires
\begin{equation*}
  8\cdot5^{j-1}\equiv39\pmod{109}.
\end{equation*}
The inverse of $8$ modulo $109$ is $41$, and
\begin{equation*}
  39\cdot41=1599\equiv73\pmod{109},
\end{equation*}
so
\begin{equation*}
  5^{j-1}\equiv73\pmod{109}.
\end{equation*}
The full family calculation combines this with the remaining branch
congruences and produces the residue class
\begin{equation*}
  j-1\equiv923\pmod{1728}.
\end{equation*}
The consistency of the representative $923$ is checked by
\begin{equation*}
  5^{923}\equiv73\pmod{109},
  \qquad
  8\cdot73=584\equiv39\pmod{109}.
\end{equation*}
For this family the first step of weight at least $3$ occurs at depth
$j+4$ and has weight $6$.  The finite base says that the first big
step of the orbit of $7$ is at depth $15$ with weight exactly $3$, so
the family is excluded.

\subsection{$d=3$ family record}

\begin{center}
\small
\begin{tabular}{@{}p{0.42\textwidth}p{0.52\textwidth}@{}}
\toprule
Item & Value \\
\midrule
surviving parameters &
  $t_j=2$, $\delta=1$, $e=2$, $u_1=u_2=1$, $a=0$ \\
state formula & $g=(8\cdot5^{j-1}-39)/109$ \\
residue & $j-1\equiv923\pmod{1728}$ \\
first big step & depth $j+4$, weight $6$ \\
contradiction & finite base: first big step at depth $15$ has weight $3$ \\
\bottomrule
\end{tabular}
\end{center}

\subsection{Full $d\ge4$ branch analysis}

For $d\ge4$, the first big step of the candidate block is one of three
types.

\begin{center}
\small
\begin{tabular}{@{}p{0.22\textwidth}p{0.40\textwidth}p{0.32\textwidth}@{}}
\toprule
Branch & First big step & Contradiction \\
\midrule
$e\ge3$ & $g_{\mathrm{prev}}\to g$, the incoming step &
  finite base forces $e=3$, $j=17$; residue modulo $640$ or $1280$
  fails \\
$e=2$, $a\ge1$ & $y\to x$, weight $1+4a\ge5$ &
  the first big step of the prefix has weight exactly $3$ \\
$e=2$, $a=0$ & $z\to w$, weight $5$ or $6$ &
  the first big step of the prefix has weight exactly $3$ \\
\bottomrule
\end{tabular}
\end{center}

In the $e\ge3$ branch, the incoming step is already big.  The finite
base allows only $e=3$ and $j=17$, and in that case the candidate
residue is excluded modulo $640$ or $1280$.  The two compiled
statements are
\texttt{dge4\_\allowbreak e3\_\allowbreak j17\_\allowbreak t1\_\allowbreak
excluded} and
\texttt{dge4\_\allowbreak e3\_\allowbreak j17\_\allowbreak
t2\_\allowbreak delta3\_\allowbreak excluded}.

In the $e=2$ branches, the first big step occurs inside the segment
with weight at least $5$.  The finite base says the first big step of
the orbit of $7$ is unique, occurs at depth $15$, and has weight
exactly $3$.  A candidate first big step of weight $5$ or $6$ cannot
occur at depth $15$ and cannot occur before it; if it occurred later,
the true depth-$15$ step would already be first.  This is recorded by
\texttt{candidate\_\allowbreak first\_\allowbreak big\_\allowbreak
step\_\allowbreak weight\_\allowbreak ge\_\allowbreak five}.

\subsection{$d\ge4$ branch table}

\begin{center}
\small
\begin{tabular}{@{}p{0.24\textwidth}p{0.36\textwidth}p{0.34\textwidth}@{}}
\toprule
Branch & First big step & Contradiction \\
\midrule
$e\ge3$ & $g_{\mathrm{prev}}\to g$ &
  finite base forces $e=3$, $j=17$; residue modulo $640$ or $1280$
  fails \\
$e=2$, $a\ge1$ & $y\to x$, weight $1+4a\ge5$ &
  the first big step of the prefix has weight exactly $3$ \\
$e=2$, $a=0$ & $z\to w$, weight $5$ or $6$ &
  the first big step of the prefix has weight exactly $3$ \\
\bottomrule
\end{tabular}
\end{center}

For the $e=2$ branches, the contradiction is the same finite fact:
the first big step of the full orbit is unique, occurs at depth $15$,
and has weight exactly $3$.  A candidate first big step of weight $5$
or $6$ therefore cannot occur at depth $15$ and cannot occur before
it; if it occurred later, the depth-$15$ step would already be the
first big step.  The $e\ge3$ branch is instead reduced by the finite
base to $e=3$, $j=17$ and excluded by the stated residues.

The $d=3$ residue class is consistent with the denominator in
$g=(8\cdot5^{j-1}-39)/109$: since
\begin{equation*}
  5^{923}\equiv73\pmod{109},
\qquad
  8\cdot73\equiv39\pmod{109},
\end{equation*}
the integer $8\cdot5^{923}-39$ is divisible by $109$.  The derivation
of the full modulus $1728$ is part of the $d=3$ family calculation.

\subsection{CRT and terminal residue}

Let $n=s-j$ and $\Delta=W_s-W_j$.  The CRT candidate $r_{j,0}$ is the
least nonnegative solution of
\begin{equation*}
  3\cdot5^n\,r_{j,0}+3B+2^\Delta
    =2^{\Delta+L+4}\left(u+m'\,2^{H_s-1}\right),
\end{equation*}
with $u$ the least nonnegative residue of
$-5^{-(L+3)}$ modulo $2^{H_s-1}$.  The terminal residue is
\begin{equation*}
  y^*=-\left(w_{\mathrm{term}}+5^{-(L+3)}\right)
  \left(3\cdot5^s\right)^{-1}\pmod{2^{H_s-1}},
\end{equation*}
and the final inequality is equivalent to $y^*\ne0$.

The Lean names are \texttt{rj0\_spec\_eq},
\texttt{terminal\_bound\_iff\_yStar\_ne\_zero}, and
\texttt{rj0\_ge\_of\_size\_conditions\_no\_hge}.

\subsection{Finite base}

The first 17 states are explicitly expanded:
\begin{equation*}
  7,9,23,29,73,183,229,573,1433,3583,4479,5599,6999,8749,21873,54683,34177.
\end{equation*}
The first big step is at depth $15$ with weight $3$.

Status: math-proven and Lean-compiled in
\texttt{FinitePrefix.lean}, with no \texttt{native\_decide}.

\subsection{Finite-prefix facts used by the exclusions}

The exclusions use the following finite facts:
\begin{itemize}[leftmargin=*]
\item the first 17 states are expanded exactly;
\item the step weights for $n<16$ are expanded exactly;
\item the first step of weight at least $3$ is at depth $15$ and has
  weight exactly $3$;
\item the first big step is unique inside this prefix;
\item a candidate first big step of weight $6$ contradicts the exact
  weight $3$ at depth $15$;
\item a candidate first big step of weight at least $5$ cannot be the
  first big step of the prefix.
\end{itemize}

The corresponding Lean names are listed in the Lean theorem index:
\texttt{fullOrbitIter\_\allowbreak prefix\_expand},
\texttt{fullOrbit\_\allowbreak prefix\_step\_weights},
\texttt{fullOrbit\_\allowbreak first\_t\_ge3\_is\_exactly\_3},
\texttt{first\_\allowbreak big\_step\_unique},
\texttt{candidate\_\allowbreak first\_big\_step\_weight\_ne\_three}, and
\texttt{candidate\_\allowbreak first\_big\_step\_weight\_ge\_five}.

\subsection{Final valuation inequality}

\begin{theorem}
Under the no-$H_{\ge}$ premises, $\FullOrbit(r)$, and $2\le H_s$,
\begin{equation*}
  \vtwo{5^{L+3}w_{\mathrm{term}}+1}\le H_s-2,
\end{equation*}
where $w_{\mathrm{term}}=(3r_s+1)/2^{L+4}$.
\end{theorem}

Status: math-proven; Lean pending in
\texttt{UnifiedCoreAudit.unified\_core\_final\_no\_hge}.
